\section*{अध्याय \ref{ch:b3:complex-spectra-magnetism} --- संकुलों के इलेक्ट्रॉनिक स्पेक्ट्रम और चुंबकत्व}

\begin{solution}{exo:b3:complex-spectra-magnetism:1}
D: E$_g$ + T$_{2g}$ ($2 + 3 = 5$); F: A$_{2g}$ + T$_{1g}$ + T$_{2g}$ ($1 + 3 + 3 = 7$); G: A$_{1g}$ + E$_g$ + T$_{1g}$ + T$_{2g}$
($1 + 2 + 3 + 3 = 9$)।
\end{solution}

\begin{solution}{exo:b3:complex-spectra-magnetism:2}
$d^1$ $^2$D $\to$ $^2$T$_{2g}$; $d^2$ $^3$F $\to$ $^3$T$_{1g}$; $d^3$ $^4$F $\to$ $^4$A$_{2g}$; $d^4$ $^5$D $\to$ $^5$E$_g$; $d^5$ $^6$S $\to$ $^6$A$_{1g}$; $d^6$ $^5$D
$\to$ $^5$T$_{2g}$; $d^7$ $^4$F $\to$ $^4$T$_{1g}$; $d^8$ $^3$F $\to$ $^3$A$_{2g}$; $d^9$ $^2$D $\to$ $^2$E$_g$।
\end{solution}

\begin{solution}{exo:b3:complex-spectra-magnetism:3}
उच्च-चक्रण \ce{Mn^{2+}} ($n = 5$) $5.92\,\mu_B$, \ce{Fe^{2+}} (4) $4.90\,\mu_B$, \ce{Co^{2+}} (3) $3.87\,\mu_B$; निम्न-चक्रण \ce{Fe^{2+}} और
\ce{Co^{3+}} ($t_{2g}^6$): 0, प्रतिचुंबकीय।
\end{solution}

\begin{solution}{exo:b3:complex-spectra-magnetism:4}
$\ce{[Mn(H2O)6]^{2+}}$ (चक्रण- और लापोर्ट-निषिद्ध, लगभग रंगहीन) $<$ $\ce{[Co(H2O)6]^{2+}}$
(लापोर्ट-निषिद्ध, कंपन-इलेक्ट्रॉनिक) $<$ $\ce{[CoCl4]^{2-}}$ (कोई सममिति केंद्र नहीं) $<$ \ce{MnO4-} (अनुमत
आवेश स्थानांतरण)।
\end{solution}

\begin{solution}{exo:b3:complex-spectra-magnetism:5}
8500: $^3$A$_{2g}\to{}^3$T$_{2g}$, अतः $\Delta_{\mathrm o} = \qty{8500}{cm^{-1}}$; 14100: $^3$T$_{1g}$(F); 24900: $^3$T$_{1g}$(P)।
$15B = 39\,000 - 3 \times 8500$, $B = \qty{900}{cm^{-1}}$, $\beta = 900/1056 = 0.85$।
\end{solution}

\begin{solution}{exo:b3:complex-spectra-magnetism:6}
$\Delta_{\mathrm o} = \qty{17400}{cm^{-1}}$; $7.5B + 1.5\Delta_{\mathrm o} - \frac12\sqrt{225B^2 - 18B\Delta_{\mathrm o} + \Delta_{\mathrm o}^2} = 24\,600$ को हल करने पर: $B = \qty{729}{cm^{-1}}$,
$\beta = 0.79$। $\nu_3 = \qty{38500}{cm^{-1}}$ (\qty{260}{nm}), पराबैंगनी में।
\end{solution}

\begin{solution}{exo:b3:complex-spectra-magnetism:7}
रत्नों में $d$ इलेक्ट्रॉन जल अणुओं की तुलना में ऑक्साइड आयनों पर अधिक
विस्थानीकृत होते हैं: वहाँ Cr--O आबंध अधिक सहसंयोजी हैं।
\end{solution}

\begin{solution}{exo:b3:complex-spectra-magnetism:8}
प्रबल ($e_g$ असमान रूप से भरा): $d^9$, उच्च-चक्रण $d^4$, निम्न-चक्रण $d^7$। शून्य: $d^3$ ($t_{2g}^3$, A पद) और
निम्न-चक्रण $d^6$ ($t_{2g}^6$)। दुर्बल: उच्च-चक्रण $d^6$ ($t_{2g}^4e_g^2$, T पद)।
\end{solution}

\begin{solution}{exo:b3:complex-spectra-magnetism:9}
अष्टफलकीय उच्च-चक्रण $d^7$: मूल पद $^4$T$_{1g}$, \omterm{def:b3:complex-spectra-magnetism:moment}{कक्षकीय योगदान} वाला एक T पद,
आघूर्ण केवल-चक्रण $3.87\,\mu_B$ से काफ़ी ऊपर। चतुष्फलकीय $d^7$ ($\equiv$ अष्टफलकीय $d^3$): मूल
पद $^4$A$_2$, कक्षकीय आघूर्ण शमित, केवल-चक्रण के निकट (उत्तेजित T पदों के साथ मिश्रण से
थोड़ा ऊपर)।
\end{solution}

\begin{solution}{exo:b3:complex-spectra-magnetism:10}
$\mu_{\mathrm{eff}} = 2.828\sqrt{1.87} = 3.87\,\mu_B$: तीन अयुग्मित इलेक्ट्रॉन ($S = \frac32$)।
\end{solution}

\begin{solution}{exo:b3:complex-spectra-magnetism:11}
$\Delta\chi_v = 3 \times 25.0/\num{400e6} = \num{1.88e-7}$; $\chi_{\mathrm m} = \num{1.88e-7}/\qty{10.0}{mol.m^{-3}} = \qty{1.88e-8}{m^3.mol^{-1}}$, अर्थात्
\qty{1.49e-3}{cm^3.mol^{-1}}; $\chi_{\mathrm m}T = 0.445$, $\mu_{\mathrm{eff}} = 1.89\,\mu_B$: एक अयुग्मित इलेक्ट्रॉन।
\end{solution}

\begin{solution}{exo:b3:complex-spectra-magnetism:12}
$T_{1/2} = 15\,000/75 = \qty{200}{K}$। \qty{250}{K} पर: $\Delta G = 15\,000 - 250 \times 75 = \qty{-3750}{J/mol}$, $K = \eu^{1.80} = 6.07$,
$x_{\mathrm{HS}} = 0.86$; $\chi_{\mathrm m}T = 0.86 \times 3.00 = \qty{2.6}{cm^3.K.mol^{-1}}$।
\end{solution}

\begin{solution}{pb:b3:complex-spectra-magnetism:1}
\textbf{1.} [Ar]$3d^3$।
\textbf{2.} $^4$F।
\textbf{3.} $^4$A$_{2g}$ + $^4$T$_{2g}$ + $^4$T$_{1g}$।
\textbf{4.} $^4$A$_{2g}$, $t_{2g}^3$ का पद, तीनों इलेक्ट्रॉन निचले कक्षकों में।
\textbf{5.} $^4$T$_{1g}$(P)।
\textbf{6.} $^4$A$_{2g}\to{}^4$T$_{2g}$, $\to{}^4$T$_{1g}$(F), $\to{}^4$T$_{1g}$(P)।
\textbf{7.} 561 और \qty{414}{nm}।
\textbf{8.} 597 और \qty{434}{nm}।
\textbf{9.} माणिक पीला-हरा और बैंगनी अवशोषित करता है और लाल (थोड़े नीले के साथ) पारगत करता है;
पन्ने के बैंड, लाल की ओर खिसके हुए, नारंगी-लाल भी अवशोषित करते हैं, और हरा
पारगत होता है।
\textbf{10.} माणिक \qty{17820}{cm^{-1}}, पन्ना \qty{16740}{cm^{-1}}।
\textbf{11.} $^4$T$_{2g}$ ($t_{2g}^2e_g$) और $^4$A$_{2g}$ ($t_{2g}^3$) की इलेक्ट्रॉन-प्रतिकर्षण ऊर्जा समान है; उनका
अंतर एक-इलेक्ट्रॉन उन्नयन ऊर्जा $\Delta_{\mathrm o}$ है।
\textbf{12.} $B = (14\,305 - 547)/15 = \qty{917}{cm^{-1}}$।
\textbf{13.} \qty{614}{cm^{-1}}।
\textbf{14.} \qty{618}{cm^{-1}}।
\textbf{15.} दोनों के लिए 0.67।
\textbf{16.} 29.0 और 27.1।
\textbf{17.} माणिक के लिए $\nu_3 = \qty{38500}{cm^{-1}}$ (\qty{260}{nm}), पराबैंगनी में, जहाँ आवेश-स्थानांतरण
अवशोषण इसे छिपा देता है।
\textbf{18.} बेरिल का अष्टफलकीय स्थल बड़ा है: लंबे Cr--O आबंध छोटा
विपाटन देते हैं, जो दूरी के साथ तेज़ी से बदलता है।
\textbf{19.} $10^7/695 = \qty{14390}{cm^{-1}}$, $23.4B$।
\textbf{20.} $^2$E$_g$ मुख्यतः $C$ पर निर्भर करता है: ऊँचा मापा गया स्तर बताता है कि $C/B$ बड़ा है,
उसी विकर्णीकरण से लगभग 5.3।
\textbf{21.} $^2$E$_g$ और $^4$A$_{2g}$ दोनों $t_{2g}^3$ के हैं, केवल एक चक्रण-उलटाव से भिन्न: आबंधन,
अतः ज्यामिति, दोनों में समान है, और संक्रमण में कोई कंपनिक
श्रेणी नहीं होती (एक 0--0 रेखा)।
\textbf{22.} यह चक्रण-निषिद्ध है।
\textbf{23.} इसकी ऊर्जा $B$ और $C$ पर निर्भर करती है, $\Delta_{\mathrm o}$ पर लगभग नहीं (आरेख पर एक सपाट रेखा),
और $B$ दोनों रत्नों में लगभग समान है।
\textbf{24.} चौड़े बैंडों में पंप किए गए आयन विश्रांत होकर दीर्घजीवी
$^2$E$_g$ स्तर में संचित होते हैं, जिससे उसमें मूल स्तर की तुलना में अधिक आयन हो सकते हैं: एक \omterm{def:b3:partition-functions:boltzmann}{समष्टि}
व्युत्क्रमण।
\textbf{25.} \textbf{$\Delta_{\mathrm o}(\text{माणिक}) - \Delta_{\mathrm o}(\text{पन्ना}) = \qty{1080}{cm^{-1}}$}: लाल और
हरे के बीच का पूरा अंतर।
\end{solution}
